\subsection{A finite known query list}
\label{sec:finite-query-list}

\begin{theorem}[A finite known query list can move evaluation to prefill]
\label{thm:finite-query-list}
Fix a cache and a list of $M$ permitted queries, known before
preprocessing. Query $q$ has a categorical output law
$p_{q,1},\ldots,p_{q,V}$. Suppose a deterministic certified algorithm,
given $q$ and an integer $t\ge1$, returns dyadic intervals of width
at most $2^{-t}$ containing all $V$ probabilities, with endpoints
of length $O(B+t)$, in at most
\[
 H(B+t)^d
\]
bit operations. Here $H\ge V$, $B\ge2$, and the fixed integer $d\ge2$
are known bounds; $B$ includes addresses and parameter lengths.
No exact real probabilities need be stored.

There is a preprocessor using
\[
 O\!\left(MH[B+\log(VH+2)]^d\right)
\]
bit operations and
$O(MV[B+\log(VH+2)])$ additional bits of storage, such that, when the
online query is supplied by its list identifier, one exact output
sample takes
\[
 O\!\left([B+\log(VH+2)]^d\right)
\]
expected bit operations. It halts almost surely, and the bound is
uniform over the query list. The cache and the certified evaluation
procedure remain available on a rare fallback branch. If $H$ is
polynomial in $T,n,D,V$ and the model encoding size, and $B$ is
polylogarithmic in those quantities, this online bound is
polylogarithmic in those quantities, regardless of the attention
temperature. Identifying a query supplied as an arbitrary vector is
a separate cost and is not covered by the list-identifier interface.
\end{theorem}
\paragraph{Proof idea.}
Prefill stores dyadic masses lying below the desired probabilities.  Their missing mass is positive but small.  Most samples finish from the stored table, while the residual branch recomputes enough certified digits to sample its exact conditional law.

\begin{proof}
The case $V=1$ is immediate. Choose a dyadic $\delta$ with
$1/(8H)<\delta\le1/(4H)$. For each query, compute dyadic lower
bounds $\ell_k\le p_k$ with
$0\le p_k-\ell_k\le\delta/V$. Such bounds follow by evaluating at
precision $t_0=\lceil\log_2(V/\delta)\rceil+2$ and rounding
downward to a common dyadic grid, clamping negative lower bounds to
zero. Put
\[
 a_k=(1-\delta)\ell_k,\qquad r=1-\sum_{k=1}^V a_k.
\]
All these numbers are stored dyadics of length
$O(B+\log(VH+2))$. They satisfy $0\le a_k\le p_k$ and
\[
 \delta\le r\le1-(1-\delta)^2\le2\delta.
\]
The left inequality uses $\sum\ell_k\le1$; the right uses
$\sum\ell_k\ge1-\delta$.

Store the prefix sums of $a_1,\ldots,a_V,r$ on their common dyadic
grid. Draw from these $V+1$ known masses using uniform grid bits
and binary search of the prefix sums. This costs a polynomial in
$B+\log(VH+2)$, covered by its $d$th power. On outcome $k\le V$,
return token $k$. On the residual outcome, use fresh randomness to
sample the categorical law
\[
 b_k=\frac{p_k-a_k}{r}.
\]
It is a probability distribution because its entries are nonnegative
and sum to one. The final probability of token $k$ is exactly
$a_k+rb_k=p_k$.

We give the precision and stopping bound for the residual sampler.
Let $U$ be a fresh uniform number on $[0,1]$, exposed through its
binary digits, and use the cumulative masses of $b$. At stage
$j\ge0$, evaluate all $p_k$ with width at most $2^{-t_j}$, where
\[
 t_j=\left\lceil\log_2(16V^2/\delta)\right\rceil+j.
\]
Each resulting residual cumulative interval has width at most
$V2^{-t_j}/r\le2^{-j}/(16V)$. Expose enough bits of $U$ that its
dyadic interval has width at most $2^{-j}/(16V)$. If this interval
is certified to lie between two adjacent cumulative boundaries,
return the corresponding token; otherwise refine.
Division by the known positive dyadic $r$ can be implemented by
cross-multiplication in every comparison, using finite integers.

Failure to decide at stage $j$ places $U$ within distance at most
$2^{-j}/(8V)$ of one of the $V-1$ true boundaries. A union bound
therefore gives probability at most $2^{-j}/4$. The probability of
never deciding is zero. The same estimate gives conditional expected
fallback work
\[
 O\!\left(
 H\sum_{j\ge0}2^{-j}[B+\log(V/\delta+2)+j]^d
 \right)
 =O\!\left(H[B+\log(VH+2)]^d\right).
\]
The sum also covers the fresh uniform bits and the $V$ cumulative
interval additions, since $H\ge V$ and $d\ge2$. Boundary equalities
have probability zero and do not require an equality oracle.

The residual branch has probability $r\le2\delta\le1/(2H)$.
Multiplying its conditional expected work by this known probability
gives the stated online bound. The initial dyadic draw contributes
the same order or less. Preprocessing performs one precision-$t_0$
evaluation and stores $V+1$ prefix masses for each of the $M$ known
queries, giving the displayed preprocessing and storage bounds.
Every arithmetic operation uses finite integers or certified dyadic
intervals.
\end{proof}

For decoding, this theorem distinguishes a finite advertised set of
next queries from an unrestricted fresh query. A finite vocabulary
can sometimes determine such a list once the previous state is fixed;
constructing that list and all its output laws is then part of
prefill. The lower bound of Section~\ref{sec:indexedlower} concerns
arbitrary fresh queries and does not apply to such a list. The tables
are specific to one cache, and a sampled token that changes the cache
requires new tables. For fixed-rank key maps,
Theorem~\ref{thm:momentdecoder} provides a decoder construction that
pays for those updates.
